\hypertarget{appendix-d8a7ddb675eb}{}
\section*{Jacobians, pushforwards, and pullbacks}
\section{Jacobians, pushforwards, and pullbacks}\label{jacobians-pushforwards-and-pullbacks}

\textbf{Curriculum depth:} Developed · \textbf{First seen:} Phase 9 · \textbf{Latest development:} Phase 9

\subsection{Definition and mechanism}\label{definition-and-mechanism}

The Jacobian is the matrix of first derivatives of a vector-valued function. It maps a small input displacement to a first-order output displacement, a forward sensitivity often called a pushforward or Jacobian-vector product. Gradients transform in the reverse direction by the transpose Jacobian, the pullback used in reverse-mode differentiation. These operations differ because vectors and covectors have different transformation rules. Jacobians compose by matrix multiplication and their determinants describe local volume changes. A Jacobian is local: it does not replace a nonlinear transformation globally. At singular points it can lose rank or become ill-conditioned, affecting inverse sensitivity, uncertainty propagation, and optimizer steps. State coordinate ordering and column-vector conventions before interpreting signs or multiplication order.

\subsection{Prerequisites}\label{prerequisites}

\begin{itemize}
\tightlist
\item
  \href{}{Derivative as a local linear map (02-Foundations-II, appendix-7bc181935c42)}
\item
  \href{}{Matrix multiplication and composition order (02-Foundations-II, appendix-ab31ef447dbf)}
\item
  \href{}{Probability density and change of measure (03-Visual-Computing-I, appendix-4815cd87ab8b)}
\end{itemize}

\subsection{Used by}\label{used-by}

\begin{itemize}
\tightlist
\item
  \href{}{Automatic differentiation and backpropagation (03-Visual-Computing-I, appendix-40df0855b4a7)}
\item
  \hyperref[appendix-6ac071f8b50c]{Hessians and local curvature}
\item
  \href{}{Nonlinear least squares: Gauss-Newton and Levenberg-Marquardt (02-Foundations-II, appendix-2c2df8345991)}
\end{itemize}

\subsection{Full teaching explanation}\label{full-teaching-explanation}

\begin{itemize}
\tightlist
\item
  \href{Continuous\%20Notes/Computing\%20Foundations\%20-\%20From\%20Bits\%20to\%20C++\%20and\%20Systems\%20APIs\#Jacobians,\%20curvature,\%20and\%20local\%20approximation}{Phase 9 teaching section}
\item
  \hyperref[appendix-adf48f5dc5c1]{Visual Recall/09 - Calculus and Optimization - Visual Recall.md} · visual recall
\end{itemize}

\subsection{Source coverage}\label{source-coverage}

Supplied lesson parts: 822 (T04 · M050), 823 (T04 · M050), 824 (T04 · M050), 825 (T04 · M050).
